\section{Introduction}

Modern automatic speech recognition has made fully local voice input practical on hardware that
users already own. This changes the design space for desktop speech interfaces. A dictation
system no longer needs to send microphone audio to a remote service in order to achieve useful
accuracy, and local processing can provide a strong privacy boundary: microphone audio can be
captured, decoded, and discarded entirely on the user's machine.

\YazSes{} was developed around this premise. The first version of our evaluation
\citep{yazses2026v1} demonstrated an offline hold-to-talk voice-input architecture on a commodity
Linux laptop and measured three Whisper checkpoints on a deterministic subset of LibriSpeech
\texttt{test-clean} \citep{panayotov2015librispeech}. That study established the feasibility of
local CPU dictation, but it also left several important questions unresolved. The evaluation used
one primary machine and operating system, focused on clean read English, treated streaming as
exploratory, and did not evaluate real meeting diarization or controlled human productivity.

The larger question is whether conclusions obtained from a single clean-speech benchmark remain
valid when an interactive system is exposed to harder audio, different recognizers, repeated
inference, different processor architectures, and longer task-specific workloads.

Our expanded measurements show that several apparently straightforward conclusions do not survive
this broader evaluation.

First, clean-speech accuracy does not fully describe robustness. In a common CPU harness,
Parakeet TDT 0.6B v2 produced a 2.06\% \wer{} point estimate on the measured clean subset and
2.88\% on the harder subset, a relatively small degradation. By comparison, Whisper
\texttt{base.en} moved from 4.01\% to 9.46\%, and Moonshine/base from 3.17\% to 8.04\%.
These data do not establish a universal model ranking, but they demonstrate that model trade-offs
inferred from clean speech alone can change substantially under harder conditions.

Second, a single aggregate \wer{} can hide an interaction failure that matters disproportionately
in dictation. Repeated \texttt{large-v3} decodes of the same hard-speech subset produced
materially different \wer{} values. Yet substitutions, deletions, and correct-word hits were
unchanged across the repeated runs used in the error decomposition. The movement came from
insertion count, which ranged from 101 to 184. For an offline benchmark, this appears as movement
in aggregate \wer{}. For interactive dictation, it is a different class of failure: a fluent
continuation that the user did not speak can be substantially more disruptive than a visibly
missing word.

Third, decoder configuration does not transfer monotonically across model sizes. Disabling
previous-text conditioning slightly worsened the measured results for \texttt{base.en}, had a
smaller negative effect for \texttt{small.en}, produced byte-identical output for
\texttt{medium.en}, and removed the observed runaway repetition mode for \texttt{large-v3}. The
average \texttt{large-v3} \wer{} difference is not, by itself, a statistically established
corpus-wide improvement: most of the aggregate gain is concentrated in a small number of
utterances. The engineering value of the configuration change is therefore better understood as
tail-risk control than as an average-accuracy optimization.

Fourth, system-level intuitions also fail when measured directly. Prepending synthetic silence
after activation cannot reconstruct speech that was never captured. Streaming partial decoding
can reveal text before release when the selected model is fast enough, but can also compete with
final decoding for CPU time and increase end-to-final latency. A simple centroid-similarity
heuristic for merging diarization clusters cannot safely repair enough split speakers without
also merging different people.

Finally, the meeting workload shows why task-specific evaluation is necessary. The original
clustering operating point produced severe over-segmentation on the complete 16-recording AMI
test split, with 75.21\% mean \der{} and an average speaker-count error of more than 155
additional clusters in four-person meetings. A meeting-specific threshold reduced mean \der{} to
26.71\%. Supplying the true number of speakers made the count itself reliable, but did not
establish an improvement in \der{} at the recalibrated threshold.

Together, these findings motivate a broader evaluation principle: \emph{offline interactive
speech systems should be evaluated as systems, not only as recognizers}. Accuracy remains
important, but deployment decisions can also depend on repeatability, catastrophic tail
behavior, latency under competing workloads, platform-specific numerical variation,
diarization calibration, and the ability of the benchmark infrastructure itself to preserve
contradictory results.

This paper therefore extends the \YazSes{} evaluation along five dimensions: alternative local
speech-recognition engines, harder speech, repeated decoding and decoder intervention,
cross-platform execution, and real meeting diarization. It also treats negative results as
first-class evidence. Several experiments were initiated to justify existing defaults and
instead showed that the underlying rationale was incomplete or wrong. Those findings changed the
system and are reported rather than discarded.
